\chapter{Вступ}

\section{Актуальність дослідження}

Сучасна дизайн-освіта переживає трансформацію під впливом цифрових технологій
та штучного інтелекту. Генеративні AI-інструменти, такі як Midjourney, DALL-E
та Stable Diffusion, змінюють практику дизайну, створюючи потребу в оновленні
навчальних програм.

В Україні функціонує понад 122 акредитованих освітніх програм
зі спеціальності 022 «Дизайн». Аналіз їх змісту дозволяє виявити тенденції
та прогалини у підготовці майбутніх дизайнерів.

\section{Мета і завдання дослідження}

\textbf{Мета:} проаналізувати зміст освітніх програм з дизайну в Україні
та розробити рекомендації щодо інтеграції AI-інструментів у навчальний процес.

\textbf{Завдання:}
\begin{enumerate}
    \item Зібрати та систематизувати дані про акредитовані програми з дизайну
    \item Провести статистичний аналіз структури навчальних планів
    \item Виявити тематичні кластери за допомогою NLP-методів
    \item Оцінити покриття AI та цифрових навичок
    \item Розробити пропозиції щодо модернізації навчальних програм
\end{enumerate}

\section{Об'єкт і предмет дослідження}

\textbf{Об'єкт:} система підготовки фахівців з дизайну в Україні.

\textbf{Предмет:} зміст та структура освітніх програм зі спеціальності 022 «Дизайн».

\section{Методи дослідження}

\begin{itemize}
    \item Веб-скрейпінг даних з порталу НАЗЯВО
    \item Статистичний аналіз (описова статистика, кореляційний аналіз)
    \item NLP-методи (тематичне моделювання, кластеризація)
    \item Мережевий аналіз (граф спільної появи курсів)
    \item Аналіз прогалин у навичках
\end{itemize}

\section{Наукова новизна}

Вперше проведено комплексний аналіз усіх акредитованих програм з дизайну
в Україні з використанням методів машинного навчання та NLP.

\section{Практичне значення}

Результати дослідження можуть бути використані:
\begin{itemize}
    \item Закладами вищої освіти для модернізації навчальних програм
    \item НАЗЯВО для розробки стандартів акредитації
    \item Методичними комісіями для оновлення освітніх стандартів
\end{itemize}

\section{Структура роботи}

Робота складається зі вступу, 10 розділів, висновків, списку використаних
джерел та додатків. Загальний обсяг роботи --- ХХ сторінок.
Робота містить ХХ таблиць, ХХ рисунків та ХХ додатків.

Емпіричну базу дослідження становлять 122 акредитаційних справ,
2,566 PDF-документів загальним обсягом 77.8 млн символів.
